% part: intuitionistic-logic
% chapter: propositions-as-types
% section: reduction

\documentclass[../../../include/open-logic-section]{subfiles}

\begin{document}

\olfileid{int}{pty}{red}

\olsection{Reduction}

In natural deduction !!{derivation}s, !!a{introduction} rule that is
followed by !!a{elimination} rule is a detour that can be ``reduced''.
For instance, an $\Intro{\land}$ followed by $\Elim{\land}$ ``cancel
out,'' since the conclusion of $\Elim{\land}$ is always one of the
conjuncts of the conjunction that \Intro{\land} produces, and the
premises of $\Intro{\land}$ are both conjuncts. Thus if we need
!!a{derivation} of either conjunct, we can simply use that
!!{derivation} without introducing the conjunction followed by
eliminating it. We have
\begin{gather*}
  \bottomAlignProof
  \AxiomC{}
  \RightLabel{$\delta_1$}
  \DeduceC{$!A_1$}
  \AxiomC{}
  \RightLabel{$\delta_2$}
  \DeduceC{$!A_2$}
  \RightLabel{$\Intro{\land}$}
  \BinaryInfC{$!A_1 \land !A_2$}
  \RightLabel{$\Elim{\land}_i$}
  \UnaryInfC{$!A_i$}
  \DisplayProof\bottomAlignProof
  \quad
  \redone
  \quad
  \AxiomC{}
  \RightLabel{$\delta_i$}
  \DeduceC{$!A_i$}
  \DisplayProof
\end{gather*}

We can consider a similar relation of reduction on the associated
proof terms suggested by this simplification. If $N_1$ is the proof
term for~$\delta_1$ and $N_2$ the one for~$\delta_2$, then we can say
the proof term for the proof on the left reduces (in one step) to the
proof term for the proof on the right:
\[
  \proj{i}{\pair{N_1, N_2}} \redone N_i
\]
In the typed lambda calculus, this is the \emph{beta reduction} rule for
the product type.

In natural deduction, a pair of inferences such as a \Intro{\land}
folloed by \Elim{\land}, i.e., a pair that is subject to reduction, is
called a \emph{cut}. In the typed lambda calculus the term on the left
of $\redone$ is called a \emph{redex}, and the term to the right is
called the \emph{reductum}. Unlike untyped lambda calculus, where only
$(\lambd[x][N])Q$ is considered to be redex, in the typed lambda
calculus the syntax is extended to terms involving $\pair{N}{M}$,
$\proj{i}{N}$, $\inj{i}{!A}{N}$, $\dcase{N}{x_1}{M_1}{x_2}{M_2}$, and
$\abort{!A}{N}$, with corresponding redexes.

Similarly we have reduction for disjunction:
\begin{gather*}
  \bottomAlignProof
  \AxiomC{}
  \RightLabel{$\delta$}
  \DeduceC{$!A_i$}
  \RightLabel{$\Intro{\lor}$}
  \UnaryInfC{$!A_1 \lor !A_2$}
  \AxiomC{$\Discharge{!A_1}{x_1}$}
  \RightLabel{$\delta_1$}
  \DeduceC{$!C$}
  \AxiomC{$\Discharge{!A_2}{x_2}$}
  \RightLabel{$\delta_2$}
  \DeduceC{$!C$}
  \DischargeRule{$\Elim{\lor}$}{x_1,x_2}
  \TrinaryInfC{$!C$}
  \DisplayProof\bottomAlignProof
  \quad
  \redone
  \quad
  \AxiomC{}
  \RightLabel{$\delta$}
  \DeduceC{$!A_i$}
  \RightLabel{$\delta_i$}
  \DeduceC{$!C$}
  \DisplayProof
\end{gather*}
This corresponds to a reduction on proof terms:
\begin{gather*}
  \dcase{\inj[!A_i]{i}{M}}{x_1}{N_1}{x_2}{N_2} \redone \Subst{N_i}{M}{x_i}
\end{gather*}
where $M$ is the proof term of~$\delta$, and $N_i$ is the proof term
of~$\delta_i$. This is the beta reduction rule of for sum types.
Here, $\Subst{M}{N_i}{x}$ means replacing all free occurrences of the
variable~$x$ in $M$ by~$N_i$.

The reduction of the function type corresponds to removal of
a detour of $\Intro\lif$ followed by a $\Elim\lif$.
\begin{gather*}
  \bottomAlignProof
  \AxiomC{$\Discharge{!A}{x}$}
  \RightLabel{$\delta$}
  \DeduceC{$!B$}
  \DischargeRule{$\Intro{\lif}$}{x}
  \UnaryInfC{$!A \lif !B$}
  \AxiomC{}
  \RightLabel{$\delta'$}
  \DeduceC{$!A$}
  \RightLabel{$\Elim{\lif}$}
  \BinaryInfC{$!B$}
  \DisplayProof\bottomAlignProof
  \quad
  \redone
  \quad
  \AxiomC{}
  \RightLabel{$\delta'$}
  \DeduceC{$!A$}
  \RightLabel{$\delta$}
  \DeduceC{$!B$}
  \DisplayProof
\end{gather*}
For proof terms, this amounts to ordinary beta reduction:
\begin{gather*}
  (\lambd[\typeof{x}{!A}][N]M)
  \redone \Subst{N}{M}{x}
\end{gather*}

In addition to reduction conversions on !!{proof}s, we can consider
permutation conversions. These apply to (sub)!!{proof}s that end in a
\Elim{\lor} followed by another !!{elimination} rule, e.g.:
\begin{multline*}
  \bottomAlignProof
  \AxiomC{}
  \RightLabel{$\delta$}
  \DeduceC{$!A_i$}
  \RightLabel{$\Intro{\lor}$}
  \UnaryInfC{$!A_1 \lor !A_2$}
  \AxiomC{$\Discharge{!A_1}{x_1}$}
  \RightLabel{$\delta_1$}
  \DeduceC{$!C_1 \land !C_2$}
  \AxiomC{$\Discharge{!A_2}{x_2}$}
  \RightLabel{$\delta_2$}
  \DeduceC{$!C_1 \land !C_2$}
  \DischargeRule{$\Elim{\lor}$}{x_1,x_2}
  \TrinaryInfC{$!C_1 \land !C_2$}
  \RightLabel{\Elim{\land}[i]}
  \UnaryInfC{$!C_i$}
  \DisplayProof\bottomAlignProof
  \quad
  \redone
  \\
  \AxiomC{}
  \RightLabel{$\delta$}
  \DeduceC{$!A_i$}
  \RightLabel{$\Intro{\lor}$}
  \UnaryInfC{$!A_1 \lor !A_2$}
  \AxiomC{$\Discharge{!A_1}{x_1}$}
  \RightLabel{$\delta_1$}
  \DeduceC{$!C_1 \land !C_2$}
  \RightLabel{\Elim{\land}[i]}
  \UnaryInfC{$!C_i$}
  \AxiomC{$\Discharge{!A_2}{x_2}$}
  \RightLabel{$\delta_2$}
  \DeduceC{$!C_1 \land !C_2$}
  \RightLabel{\Elim{\land}[i]}
  \UnaryInfC{$!C_i$}
  \DischargeRule{$\Elim{\lor}$}{x_1,x_2}
  \TrinaryInfC{$!C_i$}
  \DisplayProof
\end{multline*}
If the proof terms of $\delta$, $\delta_1$, and $\delta_2$ are $M$,
$N_1$, and $N_2$, respectively, the corresponding reduction rule for
proof terms, i.e., $\lambd$ terms, is:
\[
  \proj{i}{\dcase{M}{x_1}{N_1}{x_2}{N_2}} \redone 
\dcase{M}{x_1}{\proj{i}{N_1}}{x_2}{\proj{i}{N_2}}
\]

Of course we can eliminate detours not just at the ends of !!{proof}s,
but also within !!{proof}s, by applying the reduction conversion to a
sub-!!{proof} ending in a detour. Similarly, $\beta$-reduction applies
to subterms of $\lambd$-terms. If $N$ is a subterm of $M$, $N \redone
N'$, we can replace the subterm $N$ in~$M$ by $N'$ to obtain $M'$. In
that case we also write $M \redone M'$. If there is a sequence $M = M_1
\redone M_2 \redone M_2 \redone \dots \redone M_k = M'$ (including
where $k=1$, i.e., $M = M'$) we write $M \red M'$.

A !!{proof} where no conversion can be applied to any sub-!!{proof}
is called \emph{normal}. Similarly, a $\lambd$-term where no
conversion can be applied (to any sub-term) is called \emph{normal}. A
normal $\lambd$-term is also called a \emph{normal form}.

\end{document}
